% ============================================================================
% E11b & E16: Allocation Precision and Operator Sensitivity
% ============================================================================
% This file contains the LaTeX content for experiments E11b (allocation
% precision vs. overlap factor) and E16 (beta sensitivity).
% Insert this section after the "Robustness and Sensitivity Analysis"
% subsection (E10-E15) in the Evaluation section.
% ============================================================================

\subsection{Allocation Precision and Operator Sensitivity}

We conduct two additional experiments to address reviewer concerns about the
serial-overlap model and the $\beta$ knob: E11b measures how precisely each
policy's allocated bandwidth tracks the per-job attainable requirement
(allocation precision) across the overlap factor $\rho$, and E16 sweeps the
operator-facing $\beta$ parameter to verify that the scheduler's performance
is insensitive to this knob. All experiments use 5 random seeds.

\subsubsection{Allocation Precision under Overlap Variation (E11b)}

\begin{figure}[t]
\centering
\includegraphics[width=\columnwidth]{../sim-nextgen/outputs/figures/fig_e11b_overlap_waste.png}
\caption{E11b: Allocation Precision vs.\ Overlap Factor. Allocation precision
$\mathrm{Prec}=\sum_j \min(a_j, b_j^{\text{att}})/\sum_j a_j$, where $a_j$ is
the bandwidth allocated to job $j$ and $b_j^{\text{att}}$ its attainable
requirement, across overlap factors $\rho \in \{0.0, 0.3, 0.5, 0.85, 1.0\}$
at 500 and 630~Gbps (5 seeds, mean $\pm$ std). LongLiu achieves precision
$\approx 1.0$ at every operating point, whereas the WFS baseline
over-allocates bandwidth to jobs beyond their attainable requirement
(precision 0.70--0.88).}
\label{fig:e11b_precision}
\end{figure}

A reviewer concern with the serial compute-communication model is whether
LongLiu over-allocates bandwidth---reserving more than a job can physically
consume---and thereby wastes scarce capacity. We define the \emph{allocation
precision} $\mathrm{Prec}=\sum_j \min(a_j, b_j^{\text{att}})/\sum_j a_j$, which
equals 1 when every job receives at most its attainable requirement
$b_j^{\text{att}}$ (computed from the closed-form model with the given
$\rho$) and falls below 1 when the scheduler hands bandwidth to a job beyond
what it can use. Figure~\ref{fig:e11b_precision} reports precision as a
function of the overlap factor $\rho$ at 500 and 630~Gbps. LongLiu maintains
$\mathrm{Prec} \approx 1.0$ ($0.996$--$1.000$) at every measured point,
confirming that its closed-form allocation is exactly demand-tracking and
never over-allocates. In contrast, the WFS baseline yields precision
$0.70$--$0.88$, because its weight-based sharing ($1/c_i$) can assign a job
more bandwidth than its attainable requirement when the aggregate premium
demand is not saturated. Notably, LongLiu simultaneously delivers the higher
premium SLO attainment (82.5--87.5\% at 500~Gbps and 92.5--97.5\% at
630~Gbps, vs.\ WFS's 70--80\%), showing that the allocation is not merely
demand-tracking but also \emph{precisely} targeted at the premium jobs that
need protection.

\subsubsection{Operator Sensitivity to $\beta$ (E16)}

\begin{figure}[t]
\centering
\includegraphics[width=\columnwidth]{../sim-nextgen/outputs/figures/fig_e16_beta.png}
\caption{E16: Sensitivity to the operator knob $\beta$. Premium SLO attainment
rate (P-attn, left axis) and the average per-iteration SLO slowdown S-cont
(right axis) vs.\ the $\beta$ parameter at 500 and 630~Gbps, 5 seeds
(mean $\pm$ std). 500~Gbps uses solid lines and 630~Gbps dashed lines; blue
denotes P-attn and chocolate S-cont. P-attn remains within 85--95\% (500~Gbps)
and 92.5--97.5\% (630~Gbps) across the entire $\beta \in \{0.3, 0.5, 0.7,
1.0\}$ range, while S-cont rises monotonically from 0.978 to 0.999.}
\label{fig:e16_beta}
\end{figure}

The $\beta$ parameter is the most operator-facing knob of LongLiu's
deficit-to-priority mapping, and a natural robustness question is whether
performance hinges on its exact value. Figure~\ref{fig:e16_beta} sweeps
$\beta \in \{0.3, 0.5, 0.7, 1.0\}$ at 500 and 630~Gbps. Premium SLO attainment
is essentially flat: 85--95\% at 500~Gbps (peak at $\beta=1.0$) and
92.5--97.5\% at 630~Gbps across all four settings, with per-point variation
comparable to the seed-level noise. The average per-iteration SLO slowdown
S-cont (mean SAS of running jobs) improves monotonically from 0.978 to 0.999,
consistent with a higher $\beta$ pushing more deficit weight to the premium
pool. The wide plateau confirms that operators can choose $\beta$ freely
within $[0.3, 1.0]$ without compromising premium guarantees: the mechanism's
benefit is structural, not contingent on a fragile calibration.

\subsubsection{Summary}

Table~\ref{tab:e11b_e16_summary} summarizes the two experiments. LongLiu's
closed-form allocation is exactly demand-tracking (precision $\approx 1.0$)
under all overlap configurations---unlike the over-allocating WFS baseline---
and the premium SLO attainment is flat across the entire operator-facing
$\beta$ range, so no fragile calibration is required.

\begin{table}[t]
\centering
\caption{Summary of E11b and E16 (5 seeds each)}
\label{tab:e11b_e16_summary}
\begin{tabular}{lll}
\toprule
\textbf{Experiment} & \textbf{Key Finding} & \textbf{v4 Result} \\
\midrule
E11b (Precision) & v4 tracks attainable BW exactly & Prec $\approx$ 1.0 at all $\rho$ \\
E11b (Precision) & WFS over-allocates beyond need & Prec 0.70--0.88 \\
E16 ($\beta$) & P-attn flat across $\beta$ & 85--95\% (500G), 92.5--97.5\% (630G) \\
E16 ($\beta$) & S-cont improves with $\beta$ & 0.978 $\rightarrow$ 0.999 \\
\bottomrule
\end{tabular}
\end{table}
